\documentclass[12pt]{article}
\usepackage[T1]{fontenc}
\usepackage[utf8]{inputenc}
\usepackage{lmodern}
\usepackage{textcomp}
\usepackage{enumitem}
\usepackage{geometry}
\usepackage{graphicx}
\usepackage{amsmath, amssymb}
\geometry{margin=1in}
\begin{document}
\begin{center}
Mathematics - Undergraduate Level Assessment 15 \
Total Marks: 40 \
Answer all questions.
\end{center}

\section*{Multiple Choice (10 marks)}
\begin{enumerate}[label=\arabic*.]
\item The polynomial $x^4$ + 4 can be factored into linear factors in Z\_5[x]. Find this factorization.
\begin{enumerate}[label=(\alph*)]
    \item (x - 2)(x + 2)(x - 1)(x + 1)
    \item (x+$1)^4$
    \item (x-1)(x+$1)^3$
    \item (x-$1)^3(x$+1)
\end{enumerate}
\item How many structurally distinct Abelian groups have order 72?
\begin{enumerate}[label=(\alph*)]
    \item 4
    \item 6
    \item 8
    \item 9
\end{enumerate}
\item A tree is a connected graph with no cycles. How many nonisomorphic trees with 5 vertices exist?
\begin{enumerate}[label=(\alph*)]
    \item 1
    \item 2
    \item 3
    \item 4
\end{enumerate}
\item If a real number x is chosen at random in the interval [0, 3] and a real number y is chosen at random in the interval [0, 4], what is the probability that x \textless{} y?
\begin{enumerate}[label=(\alph*)]
    \item 1/2
    \item 7/12
    \item 5/8
    \item 2/3
\end{enumerate}
\item In xyz-space, what are the coordinates of the point on the plane 2 x + y + 3 z = 3 that is closest to the origin?
\begin{enumerate}[label=(\alph*)]
    \item (0, 0, 1)
    \item (3/7, 3/14, 9/14)
    \item (7/15, 8/15, 1/15)
    \item (5/6, 1/3, 1/3)
\end{enumerate}
\end{enumerate}

\section*{True / False (10 marks)}
\textit{Answer True or False}
\begin{enumerate}[label=\arabic*.]
\setcounter{enumi}{5}
\item True or False: If H is a normal subgroup of G, then aH = Ha for all a in G.
\item True or False: (g o $h)^2$ = $g^2$ o $h^2$ for every g, h in an abelian group G.
\item True or False: The cosets of the subgroup 4 Z in 2 Z are 4 Z and 2 + 4 Z.
\item True or False: The value of f(8) for the function f(x) defined by f(2 n) = $n^2$ + f[2(n - 1)] is 36.
\item True or False: There are exactly two abelian groups of order 4 up to isomorphism.
\end{enumerate}

\section*{Long Form Response (20 marks)}
\textit{Answer each question with a clear and complete explanation.}
\begin{enumerate}[label=\arabic*.]
\setcounter{enumi}{10}
\item In right triangle $ABC$ with $\angle B = 90^\circ$, we have $\sin A = 2\cos A$. What is $\tan A$?
\item Discuss the implications of an eigenvector v of an invertible matrix A. Explain why v must also be an eigenvector for the matrices 2 A, $A^2$, and $A^($-1).
\end{enumerate}

\vfill
\noindent\textit{End of Paper}
\end{document}